\section{Conclusion}
\label{sec:conclusion}

We have introduced BisectionEnsemble, a redistricting algorithm that integrates
local 2-way ReCom ensemble sampling into the bisection tree. By repositioning
ensemble sampling from the full plan level to the individual node level, we
avoid the large-arity full-plan split that can stall direct GerryChain
configurations on states with large seat counts. Every sampled ensemble node is
2-way, so the local ReCom sampler is never asked to solve a $p'>2$ split.
This guarantee applies to 2-way bisection nodes only; difficult local geometry
can still lower finite-$T$ acceptance, and $p'>2$ ApportionRegions nodes
continue to use METIS directly.

The algorithm is implemented as \texttt{SeedCompositor::BisectionEnsemble\{p,
ensemble\_steps\}} in \texttt{redist-cli}, using
\texttt{split\_subgraph\_bisection\_ensemble()} which seeds a METIS bisection,
runs \texttt{ensemble\_steps} ReCom steps via
\texttt{redist\_ensemble::recom::RecomChain}, and returns the plan at rank
$\lfloor p \times (\text{accepted\_count}-1) \rfloor$ when at least one
proposal is accepted, otherwise falling back to the METIS initialisation.

Empirically, the single-run observations in this paper show BisectionEnsemble
at $T=100$ improving Polsby-Popper compactness by 5--12\% over standard METIS
bisection and succeeding on the tested Texas ($k=38$) configuration where
direct GerryChain stalls. Runtime scales as $O(Tn\log n)$ under the stated
average-case assumptions and is parallelisable across tree nodes.

\paragraph{Scope and limitations.}
BisectionEnsemble is presented here as an algorithmic contribution with partisan
neutrality properties. It does not address the Voting Rights Act (Section~2)
dimension of redistricting: the algorithm has no minority-community-preservation
objective, and a BisectionEnsemble plan must be evaluated for VRA compliance
separately, using tools such as \texttt{redist analyze --types bloc-voting}
\citep{dellaLibera2026callais}. Practitioners deploying BisectionEnsemble in
jurisdictions with majority-minority district requirements should verify VRA
compliance of the final plan as a post-processing step.

\paragraph{Reproducibility.}
The \texttt{redist-cli} implementation used in the experiments is available
in the project repository. Exact replication requires the 2020 Census
P.L.~94-171 redistricting data and the VEST 2020 presidential precinct returns
\citep{vest2020}. Phase~2 will provide a complete replication archive including
per-run seeds, intermediate outputs, and criterion.rs benchmark results.

\paragraph{Future work.}
Several directions merit investigation. First, the ReCom chain's mixing time
on local subgraphs is unexplored: how many steps $T$ are needed for convergence
to the stationary distribution? Second, BisectionEnsemble currently applies only
to bisection nodes ($p' = 2$); extending it to $p'$-way nodes with $p' > 2$
(via nested bisections or $p'$-way recombination) would enable full ensemble
coverage of all ApportionRegions tree nodes. Third, the sensitivity of compactness
gains to the ReCom proposal distribution (e.g., biased spanning trees favouring
compact cuts) is a natural direction for further improvement. Finally, a formal
ergodicity proof for the local chain over the feasible bisection space of each
node would strengthen the theoretical foundation.
